\section{Learning Test Generators}
\label{sec:learning}
\subsection{Supervised Learning and Training Data}
Training changes the generator's behavior beyond one target or interaction. Supervised learning commonly uses code--test pairs or descriptions paired with tests. AthenaTest provides a historical example of learning unit-test generation with focal context~\cite{U11}. Such data teach framework syntax and common assertion patterns, but the training label also encodes a particular developer's testing choices. Multiple valid tests can exist for the same function, and a recorded test can be incomplete or redundant. Exact reproduction of a reference test is consequently a limited measure of learned testing ability.

Training corpora and task-specific adaptation supply complementary evidence. CAT-LM aligns code and test files during pretraining; UniTSyn uses language-server information to extract focal-code/test relationships across languages~\cite{T13,B09}. FuzzAug expands training data through coverage-guided fuzzing~\cite{T15}. exLong specializes in exception tests, while TOGLL and Tratto learn oracle construction~\cite{U16,O07,O09}. These examples distinguish learning repository correspondence, expanding observed inputs, and specializing a behavioral judgment. UniTSyn is principally a training-data resource, rather than another end-task evaluation benchmark.

Data filters determine the supervision: passing examples teach executability; discriminating examples additionally depend on the chosen faulty-program pool. Held-out projects and independently constructed faults test whether training transfers beyond that pool.

\subsection{Reinforcement Learning and Reward Meaning}
PyTester combines supervised learning and reinforcement learning for text-to-test generation, using automatic signals that include executable reference information~\cite{T01}. RLSQM learns from automatic feedback associated with unit-test quality~\cite{T12}. These approaches illustrate why the reward must be described in terms of the property being optimized. Syntax, execution success, structural coverage, maintainability, and semantic validity are related engineering concerns, but a reward for one does not directly establish improvement in another.

CURE makes the supervision distinction particularly clear. Its main co-evolutionary training setting omits ground-truth code supervision but pairs tasks with ground-truth unit tests, which help label generated programs and construct coder/tester rewards~\cite{T14}. Thus, absence of a reference implementation does not establish absence of reference supervision. Candidate programs, reference tests, and reference code must each be recorded, including changes across ablations.

Training access and inference access describe different dependencies. PyTester can generate from a description while training rewards execute a reference~\cite{T01}. Reference passing also assumes that the reference satisfies the requirement and that the test ignores incidental behavior.

\begin{table}[t]
\caption{Selected information-access configurations. Entries describe the cited configuration or stage; they are not claims about every variant or the underlying model's pretraining data.}
\label{tab:access}
\small
\begin{tabularx}{\linewidth}{@{}>{\raggedright\arraybackslash}p{.19\linewidth}YYY@{}}
\toprule
Configuration & Generation context & Additional evidence source & Consequence for interpretation \\
\midrule
TestGen-LLM~\cite{U07} & Existing tests and program context & Execution, stability, and coverage checks & Preserves observed regression behavior \\
AutoUAT / Test Flow~\cite{R03} & User stories, then scenarios and HTML & Human assessment in the industrial study & Design and implementation stages have different access \\
Nexus~\cite{O03} & Task information and agent interaction & Candidate implementation execution & Observations depend on candidate behavior \\
PyTester~\cite{T01} & Description and signature at inference & Reference execution during training & Inference access does not describe supervision needs \\
UTRL~\cite{T02} & Task-conditioned test generation & Reference and sampled programs during training & Adversarial learning retains a reference dependency \\
CURE, main training~\cite{T14} & Task-conditioned code/test generation & Ground-truth tests label generated programs & No ground-truth code does not mean no reference supervision \\
ConVerTest~\cite{O10} & Task and generated code/test candidates & Candidate execution and consensus & No pre-existing reference; candidate evidence remains \\
Requirements-oracle study~\cite{O12} & Requirements and minimal method structure & Researchers derive requirements from repair differences & Generator access differs from construction access \\
\bottomrule
\end{tabularx}
\end{table}

Table~\ref{tab:access} records stage-specific evidence, rather than assigning one black-box label per paper. Rewards should improve detection while constraining rejection of permitted behavior. Reference consistency is a practical proxy when alternative legal implementations are unavailable; the reference assumption remains explicit.

\subsection{Adversarial Learning and Self-Play}
UTRL learns a test generator through interaction with a code generator and uses reference and sampled programs in its training procedure~\cite{T02}. ATGen develops adversarial reinforcement learning for test-case generation with a changing adversarial curriculum~\cite{T03}. Sol-Ver studies self-play for solving and verifying code tasks~\cite{T04}. TCS separates stages aimed at reference-consistent tests and adversarial testing~\cite{T07}. CURE jointly evolves coder and tester behavior with reference-test-derived supervision in its main setting~\cite{T14}. Together, these approaches move test generation from imitating recorded tests toward learning from the behavior of generated programs.

Adaptive difficulty is valuable if it transfers beyond the current opponent. Evaluation should freeze the tester and use held-out tasks, independent error pools, and other code-model families. Matched-budget fixed or refreshed pools isolate interaction from additional erroneous examples. Agreement alone can preserve a shared misunderstanding; requirements, verified constraints, references, or human judgments anchor the labels.

\subsection{Learning Inputs and Learning for Downstream Use}
SAGE and LogiCase use grammar-based representations to connect input specifications to generated test inputs~\cite{T08,T09}. Their contribution concerns admissible and diverse inputs. An evaluator still needs a mechanism for deciding whether the resulting program behavior is correct. This division can be valuable because input validity has a different structure from expected-output reasoning and can sometimes be checked mechanically.

UTGen targets tests useful for debugging, whereas CodeRM trains and scales tests to support code reward modeling~\cite{T05,T06}. These objectives suggest an additional evaluation layer: the same test generator should be assessed both as a testing component and inside its intended downstream workflow. Improved code selection can arise from a better ranking signal even if some tests remain incorrect. Conversely, a semantically correct but weak test suite may have little effect on selection among already plausible candidates. Reporting both layers explains how the generated tests contribute to the final outcome.
